\subsection*{Rate guarantees for 32 and 64 users}

The primary comparison required a gap of at most 0.15 bit/s/Hz and a decision time of at most 250 ms for 32 users or 1,000 ms for 64. The same 512 independent 64-user channel realizations supplied nested 16-, 32- and 64-user subsets; the sizes are paired observations, not independent cohorts. This gave equal weight to the three single-site clustered delay line profiles (CDL-A, CDL-B and CDL-C) and the four-site urban macro profile. The deterministic scheduler, which updates the prices of congested tones and exchanges partners across feasible matchings, succeeded on 506 of 512 (98.83\%) and 504 of 512 (98.44\%) topologies. The prespecified within-profile bootstrap gave one-sided 97.5\% lower limits of 97.85\% and 97.27\%. Pooling the four profiles under a common binomial model gave Clopper--Pearson reference limits of 97.47\% and 96.94\%. Median gaps were 0.0613 and 0.0666 bit/s/Hz. Certificate-inclusive median times were 17.9 and 79.6 ms, and the 95th percentiles of individual calls were 33.2 and 140.3 ms (Fig.~\ref{fig:scale}a--c). All eight unsuccessful 64-user decisions arrived on time but exceeded the gap threshold, and all came from UMa. Supplementary Note~4, Tables~3--5 and Figs.~3 and 4 detail these comparisons.


How much of the improvement came from finding a better schedule? Zero-price matching, matching with power refinement and the subgradient baseline delivered identical lower rates on every topology. Those median rates were 15.1362, 15.2885 and 15.4007 bit/s/Hz at 16, 32 and 64 users. Deterministic search raised them by 0.0566, 0.0427 and 0.0294 bit/s/Hz. Its main gain was in the ceiling. At 32 and 64 users, zero-price matching's own ceiling certified a gap within 0.15 bit/s/Hz on none of the 512 topologies. Applying the smallest verified ceiling from the seven methods to the same zero-price action certified that gap on 395 and 426 topologies, respectively. This retrospective comparison measures rate quality and assigns no runtime to the combined calculation. For deterministic search, the median ratio of its lower rate to the common ceiling was 99.636\% and 99.587\% (Fig.~\ref{fig:phase}a and Supplementary Table~6).


A three-layer message-passing network supplied additional starting prices alongside the complete deterministic branch. It recovered all six failures at 32 users and five of eight at 64, retaining previous successes. Median gaps became 0.0549 and 0.0645 bit/s/Hz, with median times of 42.6 and 177.2 ms, respectively. Gains of 1.17 and 0.98 percentage points fell below the ten-point learning target; exact paired comparisons with the original baseline gave \(P=0.03125\) and \(P=0.0625\). Used alone, the network passed on 303 of 512 and 340 of 512 topologies.

We reran all 512 topologies per size with three repeats and rotating execution order. Three extra deterministic price iterations gave 507 and 505 passes at 32 and 64 users, with median gaps of 0.05942 and 0.06612 bit/s/Hz. Learned warm start gave 512 and 509 passes. Five and four topologies favoured learning, with none favouring the iteration-matched control (exact two-sided McNemar \(P=0.0625\) and \(0.125\)). The extra deterministic iterations added 2.59 and 6.78 ms to the baseline median times. Learning added 22.78 and 95.49 ms and raised the baseline's median lower rates by 0.00440 and 0.00110 bit/s/Hz.


A timing-only pilot selected larger deterministic budgets of 38 and 63 iterations. On the full cohort the median times were 38.39 and 176.35 ms, compared with 35.59 and 164.01 ms for learning in the same run. The longer deterministic search again passed on 507 and 505 topologies. Again five and four discordant outcomes favoured learning (\(P=0.0625\) and \(0.125\)). At 64 users, the longer deterministic search gave the smaller median gap (Fig.~\ref{fig:scale}d and Supplementary Table~7). With so few discordant outcomes, the tests have little resolution.


Component swaps also separated proposals from verification. On 32-user UMa, changing box order preserved a 0.074-bit/s/Hz median gap, whereas learned anchors, prices or actions widened it (Supplementary Table~4). Constructing and bounding the pair--tone table dominated 64-user deterministic time (median 50.9 ms).

For two-user analytic channels the mean gap was 0.23808 bit/s/Hz. The continuous optimum placed 0.23467 of this above the optimum and 0.00340 between the optimum and the achieved-rate floor. Median ceiling slack was only 0.00221 bit/s/Hz (IQR 0.00006--0.13947). At four, six and eight users, mean ceiling slack was at least 0.12743, 0.05487 and 0.03387 bit/s/Hz. Median reference widths were \(5.85\times10^{-5}\), \(7.34\times10^{-5}\) and \(1.50\times10^{-4}\) bit/s/Hz, but maxima of 0.514, 0.854 and 0.929 leave the mean split unresolved. Figure~\ref{fig:phase}b and Supplementary Table~8 give the endpoint quartiles. The 16-user comparison measures slack above the exact integer-table optimum, which is itself an upper bound on the physical optimum.


Propagation geometry also changes tightness. Across 64 guarded tiles per scene at 10 dB, median gaps at 32 and 64 users were 1.120 and 1.096 bit/s/Hz on the 60-GHz O1 street versus 0.102 and 0.095 on the 3.5-GHz ASU campus. Median squared channel overlaps were 0.9936 on O1 versus 0.4649 and 0.4578 on ASU, with Gram effective ranks near one versus 2.6 (Supplementary Note~5, Supplementary Table~9 and Supplementary Fig.~5). O1's nearly parallel vectors limit spatial separation. Experiments with transmitting and receiving apertures show how geometry changes focusing and resolvable modes~\cite{kludze2026,shenoam2026}; our independently transmitting users also occupy different scene footprints, so carrier frequency alone cannot explain the contrast.

At 32 users, median local-bound contribution \(T_4\) was 1.1133 on O1 versus 0.0074 on ASU. Refining 32, 128, 512 or all 7,936 pair--tone entries into 16 children changed no upper-table entry. The top-32 shortlist contained a median of two scheduled entries. Refining all 16 scheduled entries, with action and prices fixed, gave median \(T_4=1.113253\) and \(G=1.119582\) bit/s/Hz at the baseline depth of 16 children. At 64 and 256 children, median \(T_4\) fell to 1.080005 and 0.867040, respectively, whereas median \(G\) fell only to 1.118334 and 1.117494 bit/s/Hz. Unrefined alternatives keep the ceiling loose (Supplementary Table~10).

On all 64 nested two- and four-user O1 problems, the mean ceiling slack was 1.03877 and 0.90425 bit/s/Hz. The mean shortfall from the physical optimum to the action's rate floor was 0.13522 and 0.15193 bit/s/Hz. The continuous references place most of the small-system gap above the optimum. Across all 2,176 small analytic and O1 decisions, an independent evaluation of the stored filters placed their rates within \(1.41\times10^{-12}\) bit/s/Hz of their floors (Supplementary Note~7). The 32-user physical optimum remains unresolved.

\fig{fig4_scaling}{Certified gaps and decision times across price-search strategies.}{\textbf{a,b}, At each gap on the horizontal axis, the curves show the fraction of all 512 native topologies at 10 dB for which the scheduler achieved that gap or smaller within the deadline. Each topology contributes its largest gap and median time over three calls. The vertical line marks 0.15 bit/s/Hz. \textbf{c}, Filled points mark the median across topology-median wall times. Vertical caps mark the 95th percentile across all timed calls on a logarithmic time axis. \textbf{d}, Median gaps and times in the paired iteration- and time-budget reruns. Filled markers show the iteration run; open markers show the time-budget run. Circles and squares indicate 32 and 64 users. The time-budget control uses 38 and 63 iterations at 32 and 64 users, respectively, versus 10 and 9 for learning. These exploratory timings are separate from panels a--c.}{fig:scale}

\subsection*{The cost of additional pilots}
We compared power-refined zero-price matching with learned warm start on the same least-squares estimate, charging both for routine pilots. Without additional calibration pilots, learned warm start had a certified positive net advantage on 0 of 128 fixed-energy CDL-C, 28 of 128 scaled-energy CDL-C, 12 of 128 fixed UMa and 111 of 128 scaled UMa topologies. Sixteen extra pilot blocks exhausted scaled UMa's one-block budget; amortizing that charge across 256 and 1,024 blocks yielded advantages on 10 and 74 of 128 cases, respectively (Fig.~\ref{fig:phase}c,d). This accounting assumes stationary per-block rates (Supplementary Note~5, Supplementary Table~11 and Supplementary Fig.~6).

\fig{fig5_phase}{Achieved rates, remaining optimality uncertainty and the cost of pilots.}{\textbf{a}, Median certified lower rates and each method's median upper endpoint on 512 native topologies per size; the best upper endpoint is selected separately within each topology. \textbf{b}, Brackets for ceiling slack and shortfall from the optimum to the achieved-rate floor on analytic channels through eight users. Open circles mark mean endpoints, squares median endpoints, and pale bars span the lower endpoint's first quartile to the upper endpoint's third quartile. The separated 16-user native row instead compares with the exact upper-table optimum. Each row uses 512 topologies; Supplementary Table~8 gives both endpoint IQRs. \textbf{c}, Fractions of 128 independent simulated topologies per arm with a certified positive net gain, with exact two-sided 95\% binomial intervals. \textbf{d}, Scaled UMa fractions when additional pilots are charged over a stationary horizon; grey cells are unaffordable.}{fig:phase}

\subsection*{Receiver performance on successive recorded pilots}
Six-user Rice RENEW recordings supply early channels and later pilot observations~\cite{zafari,renew}. After noise whitening, the shared-reference-clock (HUB) capture achieved \(G\leq0.10\) bit/s/Hz within 250 ms on 1,867 of 1,873 post-calibration frames; the local-clock (LO) capture passed on 1,037 of 1,873. Median times were 2.75 and 2.71 ms, respectively (Fig.~\ref{fig:hw}a; Supplementary Note~6, Supplementary Table~12 and Supplementary Fig.~7). Denominators include detection failures.

Holding the action and filters fixed, every eligible later-pilot rate fell below its early floor: median differences were \(-2.501\) and \(-4.873\) bit/s/Hz on HUB and LO, respectively (Fig.~\ref{fig:hw}b). The disjoint long-training-sequence windows begin 16~\(\mu\)s apart. Filter design and evaluation on the same noisy early pilot create in-sample optimism; this replay cannot separate noise, synchronization and evolution. Dedicated experiments can isolate synchronization and reciprocity offsets~\cite{rogalin,vieira,ganesan}, including clock feedback switched over optical fronthaul~\cite{zhangsync2026}. Here each clock mode has only one independently documented acquisition layout.

We asked whether removing phase sensitivity could make the recorded-pilot floors more informative. Raw complex-difference calibration used the first 256 eligible frames and tested 1,614 HUB and 1,513 LO frames. Median full-vector floors \(L_F\) were 3.050 and 0.978 bit/s/Hz, versus projected floors \(L_P\) of 3.754 and 1.292. The later-rate medians were 21.937 and 18.927 bit/s/Hz; the median paired floor fractions were only 17.15\% and 6.76\%. Vector-to-projected radius ratios were 1.153 and 1.145, compared with 2.144 in the Gaussian benchmark. Median widths using the projected floor were 23.075 and 26.830 bit/s/Hz. All floors held, but LO achieved only 70.26\% joint residual coverage against a nominal 95\%: calibration failed to reach its target, and the rate floors held through slack.


Median absolute phase rotations were 0.0419 rad on HUB (IQR 0.0200--0.0712) and 0.2599 rad on LO (0.2331--0.2863), equivalent to 0.417 and 2.586 kHz, respectively, over 16~\(\mu\)s if entirely attributed to frequency offset. These diagnostics do not identify oscillator drift. Fixed-filter independent pilot noise predicts a projected maximum near 2.45 whitened units (97.5\% union bound 3.51), far below raw projected radii of 87.36 and 165.26 on HUB and LO.

Fixed-combiner SINR depends on projection magnitudes, permitting phase-invariant calibration. With the same split and rank 251, projection-magnitude residuals reduced projected radii to 59.52 and 70.30; phase-aligned vector radii were 81.10 and 111.72. Median projected floors rose to 6.028 and 4.564 bit/s/Hz (27.51\% and 24.46\% of paired later rates), with vector floors of 4.155 and 2.514. Median widths using the projected floor narrowed to 20.421 and 22.246 bit/s/Hz. Floor and interval coverage remained 100\%, while joint residual coverage was 98.08\% on HUB and 75.74\% on LO (Fig.~\ref{fig:hw}c,d; Supplementary Table~13). Phase invariance improves the useful floor, but does not restore nominal calibration on LO. Exhaustive fixed-equal-power enumeration of all 90 early-input schedules remained within median ceiling excesses of \(3.33\times10^{-5}\) and \(6.86\times10^{-5}\) bit/s/Hz. Mean per-entry difference powers were 18.40 on HUB and 121.62 on LO, falling after phase alignment to 12.25 and 43.95, respectively, still above the independent-noise prediction of two. Against mean early powers of 1,256.87 and 1,074.46, these aligned residuals are \(-20.11\) and \(-13.88\) dB: signal-dependent error or underestimated pilot-period noise could explain the excess, but these stored pairs cannot distinguish them (Supplementary Table~13).


The separate Argos spatial-overlap measurements are reported in Supplementary Table~14 and Supplementary Fig.~8. They describe user separation in a fixed array, complementing experiments that change array connectivity to measure signal gain~\cite{ni2026}. They do not measure a scheduled paired-payload transmission.

\fig{fig6_hardware}{Early-input certificates and chronological calibration on recorded pilots.}{\textbf{a}, Solid curves show the fraction of all 1,873 post-prefix frames with an early-input certificate below the indicated gap within 250 ms. Dashed curves show the early ceiling minus the exhaustive 90-schedule fixed-equal-power optimum, without a timing cutoff. All curves retain abstentions in the denominator; gaps above 0.30 lie outside the displayed range. \textbf{b}, Later-LTS rate of the stored action minus its uncalibrated early floor on 1,870 HUB and 1,769 LO eligible frames; the two LTS windows begin 16~$\mu$s apart. \textbf{c}, Test-frame coverage under raw and phase-invariant scores: projected floor, combined interval $[\max(L_P,L_F),U_F]$ and both residual events. The dotted line marks 95\%. \textbf{d}, Median projected floors and vector ceilings, with diamonds for later rates and triangles for vector floors. Phase-invariant intervals use projection-magnitude and phase-aligned vector residuals. HUB and LO have 1,614 and 1,513 test frames; Supplementary Table~13 gives paired widths and floor fractions.}{fig:hw}

\section*{Discussion}

On the simulated 16-antenna arrays, the unpriced matching often found a schedule close to optimal long before its own bound revealed this. A ceiling computed by another search could certify the same action. The argument extends to allocations in which disjoint pairs compete for distinct resources, provided that each pair has a valid local score bound. Different multiplexing orders, repeated user assignments or successive interference cancellation change that structure.

An interval can remain wide even when a receiver performs well. Its width contains ceiling slack \(U-V\), loss of the deployed decision \(V-\fix\), and conservatism in its floor \(\fix-L\). Our four-term identity identifies computational sources of width. Independent small-system evaluations distinguish the three rate quantities directly. On the nearly rank-one O1 channels, refining only the chosen entries moved the local term while competing entries kept the ceiling high. These observations point to different uses of extra computation depending on which entries control the bound.

The message-passing proposals recovered a few difficult topologies. Their small advantage in pass counts over the two deterministic controls remains uncertain, and extra deterministic iterations gave a smaller median gap at 64 users. The primary median times were 18--80 ms. Deployment would require channels that remain informative over that delay or allocation at a slower control layer. Millisecond scheduling would require a reduction in computation and a new timing evaluation.

The deployed filters also determine which channel errors matter. Calibrating projection magnitudes gave median floors of 24--28\% of the later-pilot rate, yet residual coverage fell short of 95\% in the local-clock recording. If quiet-slot variance understates active-period noise and distortion, hardware rates are optimistic and the gaps require recalculation with an appropriate noise model. A live test would need independent acquisitions, transmission of the selected pairs and powers, and decoded payload at the scheduling delay.
